\chapter{Bar duality with ordered collision coefficients}
\label{chap:ordered-coefficient-bar}

A product of two insertions must retain the coefficient on their
collision region. With three insertions, its two iterated products
must also agree on a common region of expansion. The coefficient
algebra and the state multiplication therefore enter the same
associativity equation.

The two-branch algebra of
Chapter~\ref{chap:ordered-collisions-residues} determines the first
coefficient problem. For several labels, pairwise choices of order
must come from one total order. Three labels already distinguish
these conditions: there are six total orders, while three independent
binary choices give eight possibilities. The two extra choices are
cycles. We retain all total orders, construct their common
coefficients, and place a nonzero collision class inside an
associative multiplication.

The resulting bar differential will join two states using that
class. Its dual differential will then recover the same coefficient.
This gives an explicit relation between collision geometry and
Koszul duality.

\section{Coefficients on all order charts}

Let $I$ be a finite nonempty set of labels and put
$R_I=\mathbb C[z_i\mid i\in I]$.
Each $z_i$ is the holomorphic position of the corresponding label
on a complex line. A total order $\sigma$ lists every label exactly
once. Write $\mathcal O_I$ for the finite set of these orders.
Take one coordinate chart with polynomial ring $R_I$ for each
$\sigma\in\mathcal O_I$.

For a nonempty subset $J\subset\mathcal O_I$, let $D(J)$ be the set
of unordered pairs whose relative order varies among the orders in
$J$. Their common region has coefficient ring
\[
 R_{I,J}=R_I[(z_i-z_j)^{-1}\mid\{i,j\}\in D(J)].
\]
The sign chosen for a difference does not affect its localization.
Two charts agree on this region by the identity on coordinates.
For three orders $\sigma,\tau,\rho$, the pairs varying among them
are exactly the union of the pairs varying between $\sigma,\tau$
and between $\sigma,\rho$.
Thus the common region is independent of the anchor chart.
The identity identifications consequently agree on triple overlaps.

At fixed complex positions, divide $I$ into blocks of equal positions.
Two total orders are identified precisely when their restrictions
to each equality block agree.
Every tuple of block orders extends to a total order of $I$.
The number of retained branches is therefore
\[
 \prod_B |B|!,
\]
where $B$ runs over the equality blocks.
In particular, a total three-point collision has six branches.
The binary assignments $1<2<3<1$ and $1>2>3>1$ cannot occur
in a total order.

We need coefficients that retain comparisons on intersections,
including comparisons between those comparisons.
For each $J$, introduce symbols $t_\sigma$ of degree zero and
$dt_\sigma$ of degree one, for $\sigma\in J$.
Form
\[
 \Omega_J=
 \mathbb C[t_\sigma\mid\sigma\in J]
 \otimes\Lambda(dt_\sigma\mid\sigma\in J)
 \big/\left(\sum_{\sigma\in J}t_\sigma-1,
             \sum_{\sigma\in J}dt_\sigma\right).
\]
Set $d(t_\sigma)=dt_\sigma$ and $d(dt_\sigma)=0$, and extend by
the graded Leibniz rule. The two displayed relations are preserved.
The symbols are polynomial differential forms on a simplex:
the $t_\sigma$ are its barycentric coordinates.
Only this polynomial algebra of forms is used.

For $K\subset J$, restriction to the face $K$ sets $t_\sigma$ and
$dt_\sigma$ to zero for $\sigma\notin K$.
The coefficient map $R_{I,K}\to R_{I,J}$ is localization.
Define $C_I$ to consist of the families
\[
 c=(c_J)_{\varnothing\ne J\subset\mathcal O_I},
 \qquad c_J\in R_{I,J}\otimes\Omega_J,
\]
whose restrictions to each face $K\subset J$ equal the localization
of $c_K$.
Products and the differential are taken componentwise, with
$dz_i=0$.
They preserve the face equations, so $C_I$ is a commutative
differential graded algebra.
Its role is to retain compatible coefficients and their comparison
forms on every intersection.

\section{A nonzero collision class}

Fix distinct labels $i,j$. On the component indexed by $J$, put
\[
 q_{ij,J}=\sum_{\substack{\sigma\in J\\i\prec_\sigma j}}t_\sigma.
\]
These polynomials agree on faces.
If the relative order of $i,j$ is constant on $J$, then
$q_{ij,J}$ is either zero or one, so its differential vanishes.
Otherwise $z_i-z_j$ is invertible in $R_{I,J}$.
Thus the componentwise formula
\[
 \omega_{ij}=\frac{dq_{ij}}{z_i-z_j}
\]
defines a degree-one element of $C_I$: use zero on components where
the numerator vanishes and the denominator was not inverted.
It is closed, since $d^2q_{ij}=0$ and $d(z_i-z_j)=0$.
Also $q_{ji}=1-q_{ij}$, so $\omega_{ji}=\omega_{ij}$.

\begin{proposition}\label{prop:ocb-nonzero-coefficient}
The class $[\omega_{ij}]$ in $H^1(C_I)$ is nonzero.
\end{proposition}

\begin{proof}
Choose two total orders differing only by an adjacent exchange
of $i,j$. Their intersection inverts only $z_i-z_j$.
On their two-vertex simplex, choose a coordinate $t$ increasing
from the first vertex to the second.
The restriction of $\omega_{ij}$ is
$\pm dt/(z_i-z_j)$.

If $\omega_{ij}=df$, the component of $f$ on this edge is a
polynomial in $t$ with localized coefficients.
Its values at $t=0$ and $t=1$ lie in $R_I$, by face compatibility.
Integrating its derivative from zero to one would give
\[
 f(1)-f(0)=\frac{\pm1}{z_i-z_j}.
\]
Polynomial integration is defined coefficient by coefficient over
$\mathbb C$.
The left side belongs to $R_I$.
The right side does not: multiplying by $z_i-z_j$ and then setting
$z_i=z_j$ would give $0=\pm1$.
This contradiction proves the claim.
\end{proof}

For two labels, this is precisely the principal part of the
one-coordinate collision from
Proposition~\ref{prop:ocr-line-algebra}, with the comparison written
as a polynomial one-form on an interval.

\section{An associative multiplication with a collision coefficient}

An odd coefficient such as $\omega_{ij}$ cannot multiply two
degree-zero states into another degree-zero state without a shift.
Introduce states $a_i$ of degree zero and $b_{ij}$ of degree $-1$.
Let
\[
 A_I=C_I1\oplus\bigoplus_{i\in I}C_Ia_i
           \oplus\bigoplus_{i\ne j}C_Ib_{ij}.
\]
This is a finite free graded $C_I$-module.
Set $da_i=db_{ij}=0$ and use the coefficient differential.
The coefficient action is graded central:
$vc=(-1)^{|v||c|}cv$ for homogeneous elements.
Declare $1$ to be the unit and define
\begin{equation}\label{eq:ocb-multiplication}
 a_i a_j=\omega_{ij}b_{ij}\quad(i\ne j),
 \qquad a_i^2=0.
\end{equation}
Every product involving a $b_{ij}$ and another nonunit state is zero.
Extend by bilinearity with the graded coefficient signs.
The augmentation $A_I\to C_I$ is identity on coefficients and kills
all $a_i,b_{ij}$.
Its kernel $\bar A_I$ is the augmentation ideal.

\begin{proposition}\label{prop:ocb-coefficient-algebra}
These formulas define an augmented associative differential graded
algebra over $C_I$.
\end{proposition}

\begin{proof}
Every product of three augmentation-ideal generators vanishes,
under either parenthesization.
Products involving the unit satisfy associativity by its definition.
Moving homogeneous coefficients to the left gives the same total
graded sign in both parenthesizations, so this proves associativity
for arbitrary elements.
The only nonzero state products have closed coefficient $\omega_{ij}$.
Their degrees are $0+0=1+(-1)$.
Thus their differential and the Leibniz expression both vanish.
The coefficient Leibniz rule proves the assertion for all elements.
\end{proof}

The states $b_{ij}$ and $b_{ji}$ are distinct.
Thus equality of the coefficients $\omega_{ij}=\omega_{ji}$ does
not impose commutativity of the state multiplication.
The order of the inputs remains in the output state.

\section{The ordered word differential}

To retain all iterated multiplications, suspend the augmentation
ideal downward by one degree.
Write $sv$ for an element with $|sv|=|v|-1$.
Its scalar action satisfies
$s(cv)=(-1)^{|c|}c\,sv$ for homogeneous $c$.
The word module is
\[
 B_{C_I}A_I=\bigoplus_{r\ge0}(s\bar A_I)^{\otimes_{C_I}r}.
\]
Length zero is $C_I$, with its coefficient differential.
Write $[v_1|\cdots|v_r]$ for
$sv_1\otimes\cdots\otimes sv_r$.
An element has finitely many nonzero word components.
Cutting a word at every possible gap defines its deconcatenation
coproduct, including the two empty end pieces.

The differential is determined by the internal differential
$b_1(sv)=-s(dv)$ and the adjacent multiplication
\[
 b_2(sv\otimes sw)=(-1)^{|v|}s(vw).
\]
Insert these maps at every input or adjacent pair with the graded
tensor sign. In particular, the multiplication part is
\begin{equation}\label{eq:ocb-bar-faces}
 b_{\mathrm{mult}}[v_1|\cdots|v_r]
 =\sum_{i=1}^{r-1}
 (-1)^{\sum_{j<i}(|v_j|-1)+|v_i|}
 [v_1|\cdots|v_iv_{i+1}|\cdots|v_r].
\end{equation}
It has degree one because suspension is removed once at the
contracted pair. The internal part also has degree one.

\begin{proposition}\label{prop:ocb-bar-square}
The sum $b=b_1+b_{\mathrm{mult}}$ satisfies $b^2=0$.
It commutes with deconcatenation in the graded differential sense.
\end{proposition}

\begin{proof}
Two internal differentials on different letters cancel by the tensor
signs. On one letter they give $d^2=0$.
Two adjacent contractions with disjoint input pairs also cancel:
each has degree one, so interchanging them reverses the sign.
On three overlapping inputs, direct substitution gives
\[
 b_{\mathrm{mult}}^2[v|w|u]
 =(-1)^{|w|}s\bigl((vw)u-v(wu)\bigr)=0.
\]
The remaining terms apply one internal differential and one product.
Their sum is the suspension of
$d(vw)-(dv)w-(-1)^{|v|}v\,dw$, which is zero.
Every term in $b^2$ occurs in one of these cases.

A word cut disjoint from the inputs of a differential term can be
made before or after that term.
The two choices have the same tensor sign.
When a pair is multiplied, only cuts outside that pair survive.
These are exactly the terms obtained by differentiating either
factor of the deconcatenation coproduct.
\end{proof}

For the collision multiplication~\eqref{eq:ocb-multiplication},
the first formulas are
\[
 b[a_i|a_j]=[\omega_{ij}b_{ij}],
\]
\[
 b[a_i|a_j|a_k]
 =[\omega_{ij}b_{ij}|a_k]
  -[a_i|\omega_{jk}b_{jk}].
\]
The next multiplication differential vanishes on both terms because
each contains a $b$-state.
The comparison coefficient remains a cochain throughout these
operations. It is not replaced by a scalar residue.

\section{The dual differential records the same collision}

Assign weight one to $a_i$, weight two to $b_{ij}$, and weight
zero to coefficients. The multiplication and the bar differential
preserve total weight.
Each weight contains finitely many words in the finite state basis.
It is therefore a finite free graded $C_I$-module.

Completion must retain the cochain grading. For any complex
$M=\bigoplus_{w\ge0}M(w)$ whose differential preserves weight, define
\[
 (\widehat M)^n=\prod_{w\ge0}M(w)^n,
 \qquad
 \widehat M=\bigoplus_{n\in\mathbb Z}(\widehat M)^n.
\]
Thus an element may have infinitely many weight components but only
finitely many total cochain degrees. All weight completions in this
chapter use this convention.

Dualize each weight over $C_I$ with the same degreewise product:
\[
 (A_I^!)^n=\prod_{w\ge0}
 \operatorname{Hom}_{C_I}(B_{C_I}A_I(w),C_I)^n,
 \qquad A_I^!=\bigoplus_n(A_I^!)^n.
\]
This defines the augmentation dual
\[
 A_I^!=\operatorname{Hom}_{C_I}(B_{C_I}A_I,C_I).
\]
Maps in this graded Hom obey
$F(cv)=(-1)^{|F||c|}cF(v)$.
The differential is
$dF=d_{C_I}F-(-1)^{|F|}Fb$.
Define a product of such maps by deconcatenation and multiplication
in $C_I$, using the graded tensor convention.
Coassociativity of cuts proves associativity of this product.
The differential identity follows from
Proposition~\ref{prop:ocb-bar-square} and the coefficient Leibniz rule.

Let $\alpha_i$ and $\beta_{ij}$ be the duals of $sa_i$ and
$sb_{ij}$, with values one on those basis vectors and zero on the
others. Their degrees are one and two, respectively.
Their weights are one and two.
Word duality identifies the algebra just defined with
\[
 C_I\langle\!\langle\alpha_i,\beta_{ij}\rangle\!\rangle.
\]
The double brackets denote this degreewise completion by positive
weight. Word coefficients must have finite support in total cochain
degree, including the degree of their coefficients in $C_I$.
For example, when $I$ has one label, $C_I=\mathbb C[z]$ and
$|\alpha|=1$. The series $\sum_{r\ge0}\alpha^r$ is excluded:
it would have a nonzero component in every degree.
Multiplication has finitely many decompositions in each fixed weight.

There is an additional finiteness in the present coefficient algebra.
Put $N=|\mathcal O_I|$. Simplex forms have degrees from zero to
$N-1$, and a state bar word of weight $w$ has degree $-w$.
In fixed degree, only finitely many weights therefore occur in either
the bar or its dual. Their degreewise completion adds no elements
in this model. The convention becomes essential for cobar and for
coefficient algebras with different degree bounds.

\begin{proposition}\label{prop:ocb-dual-differential}
The augmentation dual has differential
\begin{equation}\label{eq:ocb-dual-differential}
 d\alpha_i=0,\qquad
 d\beta_{ij}=-\omega_{ij}\alpha_i\alpha_j,
\end{equation}
together with the differential on $C_I$.
\end{proposition}

\begin{proof}
No product has output $a_i$, so $d\alpha_i=0$.
The only state product detected by $\beta_{ij}$ is $a_i a_j$.
Suspension gives
$s(\omega_{ij}b_{ij})=-\omega_{ij}sb_{ij}$.
Since $|\beta_{ij}|=2$, the Hom differential therefore evaluates
to $+\omega_{ij}$ on $sa_i\otimes sa_j$.
By the tensor sign,
\[
 (\alpha_i\alpha_j)(sa_i\otimes sa_j)
 =(-1)^{|\alpha_j||sa_i|}=-1.
\]
Thus the right side of~\eqref{eq:ocb-dual-differential} has exactly
the required value. It is zero on all other basis words.
This determines the differential on the generators.
\end{proof}

The geometric coefficient appears as a quadratic differential in
the dual algebra. Its presence expresses the original state
multiplication, including its cohomological degree and its order.

\section{Why the augmentation dual computes derived endomorphisms}

Denote the augmentation by $\epsilon:A_I\to C_I$.
It makes $C_I$ a left $A_I$-module.
As in Chapter~\ref{chap:relations-homotopies}, derived endomorphisms
are computed from a resolution by taking module maps and their homotopies.
An explicit resolution identifies the dual word algebra with this
endomorphism algebra up to quasi-isomorphism.

Put $P_I=A_I\otimes_{C_I}B_{C_I}A_I$.
In addition to the internal and bar differentials, include the
leading face
\[
 d_{\mathrm{lead}}\bigl(a_0\otimes[a_1|\cdots|a_r]\bigr)
 =-(-1)^{|a_0|}a_0a_1\otimes[a_2|\cdots|a_r].
\]
It is zero on length zero.
Associativity gives its cancellation with the first adjacent bar
face. Disjoint faces and internal differentials cancel by the
same signs as in Proposition~\ref{prop:ocb-bar-square}.
Thus the full differential on $P_I$ squares to zero.
Augmentation on length zero, and zero on positive lengths,
defines a map $P_I\to C_I$.

For homogeneous $a_0$, write
$\bar a_0=a_0-\epsilon(a_0)1$.
The degree-minus-one map
\[
 h\bigl(a_0\otimes[a_1|\cdots|a_r]\bigr)
 =-1\otimes[\bar a_0|a_1|\cdots|a_r]
\]
contracts this augmentation as a $C_I$-complex.
To verify the identity, insert the leading letter and expand
$dh+hd$.
Internal differentials on that letter cancel because
$d\bar a_0=\overline{da_0}$.
Internal differentials and adjacent faces to its right cancel
in pairs with their suspension signs.
The leading face contributes
$\bar a_0\otimes[a_1|\cdots|a_r]$.
When $r>0$, the next face contributes the remaining
$\epsilon(a_0)\otimes[a_1|\cdots|a_r]$.
Indeed, the two relevant signs are $(-1)^{|a_0|}$ from the face
and $(-1)^{|a_0|}$ from suspending
$\epsilon(a_0)a_1$, whose product is one.
Thus $dh+hd=1$ in positive word length.
In length zero it equals $a_0\mapsto a_0-\epsilon(a_0)1$.
This proves that $P_I\to C_I$ is a quasi-isomorphism.

As a graded left $A_I$-module, $P_I$ is free on the state words.
The differential of a basis word involves shorter words and
coefficient differentials.
Such a resolution is called semifree.
It permits a map or a homotopy to be constructed in increasing
word length. At each step its remaining equation asks for a
primitive of a cycle in the target complex.
In an acyclic target that primitive exists.
The finite state basis gives finitely many words at each length,
and every differential equation uses only earlier choices.
Consequently $\operatorname{Hom}_{A_I}(P_I,-)$ preserves
quasi-isomorphisms, by applying this construction to their cones.

Restriction to the words with leading coefficient $1$ identifies
\[
 \operatorname{Hom}_{A_I}(P_I,C_I)
 =\operatorname{Hom}_{C_I}(B_{C_I}A_I,C_I).
\]
The leading face vanishes under this Hom because every $a_1$
belongs to the augmentation ideal.
Thus the differential is exactly the dual differential already
computed. This defines a model for
$R\operatorname{Hom}_{A_I}(C_I,C_I)$.

Its product also has the endomorphism interpretation.
Deconcatenation makes $P_I$ a right bar comodule: cut the word,
retaining its left piece in $P_I$.
A homogeneous word functional $F$ lifts to an endomorphism
\[
 \widetilde F=(1\otimes F)\circ\Delta_{P_I}.
\]
The tensor rule includes the sign for moving $F$ past the left
piece. Cutting twice shows that convolution of functionals
becomes composition of these lifts.
The coaction commutes with the differentials, including the
leading face, so this statement holds for the differential algebras.
Composing a lift with $P_I\to C_I$ recovers the original functional.
The induced map
$\operatorname{End}_{A_I}(P_I)\to
\operatorname{Hom}_{A_I}(P_I,C_I)$
is a quasi-isomorphism by the semifree property.
Its section by the lifts is therefore a multiplicative
quasi-isomorphism. This proves the asserted derived endomorphism
interpretation with its product.

\section{Recovering the algebra from its bar coalgebra}

Complete the bar by positive weight and denote it by $Q_I$.
At each weight this retains the same finite free word module.
The completion is degreewise, so here $Q_I=B_{C_I}A_I$ as a
graded module. Coproducts take values in the completed tensor product
specified by
\[
 (M\widehat\otimes_{C_I}N)^n
 =\prod_{w\ge0}\left(
   \bigoplus_{a+b=w}M(a)\otimes_{C_I}N(b)
   \right)^n,
\]
with a direct sum over $n$ as above.
The reduced coproduct $\bar\Delta$ omits the two empty end cuts.
Write $t$ for suspension upward by one degree, so $|tc|=|c|+1$.
The completed cobar algebra is the tensor algebra on the elements
$tc$ with $c$ in the positive-weight part of $Q_I$, completed by
weight in each cochain degree. This completion can add elements:
$t(sa_i)$ has degree zero, so $\sum_{r\ge0}(t(sa_i))^r$ is
allowed. Its differential on a generator is
\[
 d_\Omega(tc)=-t(bc)
 +\sum_{\bar\Delta c=\sum c'\otimes c''}
       (-1)^{|c'|+1}(tc')(tc'').
\]
Extend it by the graded Leibniz rule.
The internal square vanishes by $b^2=0$.
The mixed terms cancel because $b$ commutes with the coproduct.
Two coproduct terms cancel by coassociativity and their suspension
signs. Thus $d_\Omega^2=0$.

Map $t(sa)$ to $a\in\bar A_I$, and map generators from longer
bar words to zero. Extending multiplicatively gives the counit
\[
 \widehat\Omega_{C_I}Q_I\longrightarrow A_I.
\]
On a two-letter bar generator, the internal and coproduct terms
map to $-(-1)^{|a|}ab$ and $(-1)^{|a|}ab$.
They cancel. Longer words map to zero in the same differential
identity, and single letters commute with the internal
differential. Hence the counit is a chain map.

\begin{theorem}\label{thm:ocb-reconstruction}
The displayed counit is a quasi-isomorphism in every weight,
and after the declared weight completion.
\end{theorem}

\begin{proof}
Expand a cobar word into its original state letters.
At weight $w$ there are at most $w$ letters, since each has
positive weight. Filter by this number.
The part using state multiplication lowers the filtration.
On its associated graded, each gap between consecutive letters
has two possible positions: inside one bar block, or between two
cobar factors.
The coproduct differential changes the first position to the
second with coefficient $1$ or $-1$.

After the suspension signs, this contribution is the tensor
product of one copy of $[\mathbb C\xrightarrow{1}\mathbb C]$
for each gap, tensored with the state letters and coefficients.
If a gap exists, contract the leftmost gap by the inverse map,
with the tensor sign from preceding factors.
The contractions on other gaps cancel in $dh+hd$, while the
chosen gap gives the identity.
Only the one-letter summand survives, and the counit is identity
on that summand. It is also identity on the weight-zero unit.
The finite filtration proves the weightwise quasi-isomorphism.

Finally, products of exact sequences of modules remain exact:
a preimage of a tuple in a product of surjections is chosen
component by component.
Applied in each cochain degree, this shows that products preserve
cohomology of these weightwise complexes.
The degreewise product of the proved equivalences is the completed
equivalence. The target $A_I$ has only weights zero, one, and
two, so completion leaves it unchanged.
\end{proof}

\section{Nested positions and a common expansion ring}

Place labels $1,2$ in a small cluster near $Z$, while label $3$
stays near $W$. Introduce a scale $\varepsilon$ and set
\[
 z_1=Z+\varepsilon x_1,\qquad
 z_2=Z+\varepsilon x_2,\qquad z_3=W.
\]
The inner order distinguishes $1<2$ from $2<1$.
The outer order distinguishes the cluster preceding $3$ from
the cluster following $3$.
The four resulting global orders are $123,213,312,321$.

On an intersection of these charts, invert $x_1-x_2$ if the inner
order varies and invert $Z-W$ if the outer order varies.
Let $D$ denote the resulting localization of
$\mathbb C[Z,W,x_1,x_2]$.
Expand in the ring
$D[[\varepsilon]][\varepsilon^{-1}]$.
The exponent of $\varepsilon$ in each element has a fixed lower
bound, as required by Section~\ref{sec:cc-localization}.

The inner reversed difference has inverse
\[
 (z_1-z_2)^{-1}
 =\varepsilon^{-1}(x_1-x_2)^{-1}.
\]
On an outer overlap, geometric series give
\[
 (z_1-z_3)^{-1}
 =\sum_{r\ge0}\frac{(-\varepsilon)^r x_1^r}{(Z-W)^{r+1}},
 \qquad
 (z_2-z_3)^{-1}
 =\sum_{r\ge0}\frac{(-\varepsilon)^r x_2^r}{(Z-W)^{r+1}}.
\]
Multiplying each series by its proposed denominator gives one
coefficient by coefficient. This proves that the substitutions
respect the localized rings.

Let $q_{\mathrm{in}}$ and $q_{\mathrm{out}}$ be the sums of the
barycentric coordinates recording the inner and outer orders on
this four-chart simplex. Explicitly,
\[
 q_{\mathrm{in}}=t_{123}+t_{312},\qquad
 q_{\mathrm{out}}=t_{123}+t_{213}.
\]
Pullback of simplex forms sends each
global barycentric coordinate to the sum of the coordinates of
charts inducing that global order.
It commutes with $d$ and with face restriction.
The collision coefficients consequently expand as
\[
 \omega_{12}\longmapsto
 \varepsilon^{-1}\frac{dq_{\mathrm{in}}}{x_1-x_2},
\]
\[
 \omega_{13}\longmapsto
 dq_{\mathrm{out}}\sum_{r\ge0}
 \frac{(-\varepsilon)^r x_1^r}{(Z-W)^{r+1}}.
\]
These expressions remain comparison forms of degree one.

More generally, for any map of commutative differential graded
coefficient algebras $g:C_I\to E$, define
$A_{I,E}=E\otimes_{C_I}A_I$.
Its multiplication replaces $\omega_{ij}$ by $g(\omega_{ij})$.
There is a strict isomorphism of bar complexes
\[
 E\otimes_{C_I}B_{C_I}A_I
 \cong B_E A_{I,E}.
\]
It sends every tensor word to the identical word with extended
coefficients. Formulas~\eqref{eq:ocb-bar-faces} agree term by term,
so it preserves the differential and every cut.
For the completed dual, extend coefficients in each finite weight
and then take the inverse limit in each cochain degree, retaining
a direct sum over degrees. The completed tensor convention is the
one displayed above with $E$ in place of $C_I$.
This order of operations gives the same dual differential with
$g(\omega_{ij})$. If $E$ is unbounded in degree, infinitely many
weights can contribute to one degree; the finite-degree support
condition still applies.

In the displayed three-point expansion, for example,
\[
 d\beta_{13}
 =-dq_{\mathrm{out}}\sum_{r\ge0}
 \frac{(-\varepsilon)^r x_1^r}{(Z-W)^{r+1}}
 \alpha_1\alpha_3.
\]
Thus nested collision expansion commutes with an actual bar
operation and an actual dual differential.

The state $b_{12}$ still carries two labels.
A chiral collision map must additionally produce a state at their
merged position, with compatible translation and support.
Such a map must commute with the bar faces above, or supply
homotopies for their differences satisfying the higher face
identities. Coefficient expansion provides the common domains on
which those equations can be stated.
